% GENERATED FILE. Edit canonical lessons, then run npm run generate:books.
% canonical-source-hash: fnv1a64:e2e279df89badbbd
% canonical-lessons: TE-C295-vijayam, TE-C295-akaram, TE-C295-endukaina, TE-C295-enduko, TE-C295-elagaina, TE-R295-first-pass-words-from-chapters-287-290, TE-R295-second-pass-words-from-chapters-291-295

\chapter{Success, Shape, In case, For some reason, Somehow}
\label{ch:te-success-shape-in-case-for-some-reason-somehow}
\hlchaptermodality{295}

\begin{chapteropening}
\textbf{By the end of this chapter:} \emph{I can say success, a shape, in case, for some reason and somehow.}

Complete the last lesson of chapter 295: I can say success, a shape, in case, for some reason and somehow.
\end{chapteropening}

\section[\texorpdfstring{\te{విజయం} (vijayaṁ)}{విజయం (vijayaṁ)}]{\te{విజయం} (vijayaṁ) — success}

\label{lesson:TE-C295-vijayam}



\begin{quote}
Before the new one: say the Telugu for almost, then the Telugu for all in all.
\end{quote}

\subsection*{You'll want to know: \te{విజయం}}
\textbf{\te{విజయం}} — \emph{vijayaṁ} — \textquotedblleft{}success\textquotedblright{}.

\begin{grammarlens}[title={What you've built}]
One more word for this chapter.
\end{grammarlens}

\subsection*{Your turn}
\begin{itemize}
\raggedright
  \item \emph{Say it:} \emph{vijayaṁ}
  \item \emph{Say it:} \emph{vijayaṁ}, once more
  \item \emph{Say it:} say \emph{vijayaṁ}
  \item \emph{Recall:} say the Telugu for almost, then the Telugu for all in all, then say \emph{vijayaṁ} again
\end{itemize}

\begin{tcolorbox}[breakable,colback=teal!4,colframe=teal!35!black,title={Before you move on}]
What does \te{విజయం} mean? (\textquotedblleft{}Success\textquotedblright{}.) Say it once more.
\end{tcolorbox}
\section[\texorpdfstring{\te{ఆకారం} (ākāraṁ)}{ఆకారం (ākāraṁ)}]{\te{ఆకారం} (ākāraṁ) — a shape}

\label{lesson:TE-C295-akaram}



\begin{quote}
Before the new one: say the Telugu for all in all, then the Telugu for success.
\end{quote}

\subsection*{You'll want to know: \te{ఆకారం}}
\textbf{\te{ఆకారం}} — \emph{ākāraṁ} — \textquotedblleft{}a shape\textquotedblright{}.

\begin{grammarlens}[title={What you've built}]
One more word for this chapter.
\end{grammarlens}

\subsection*{Your turn}
\begin{itemize}
\raggedright
  \item \emph{Say it:} \emph{ākāraṁ}
  \item \emph{Say it:} \emph{ākāraṁ}, once more
  \item \emph{Say it:} say \emph{ākāraṁ}
  \item \emph{Recall:} say the Telugu for all in all, then the Telugu for success, then say \emph{ākāraṁ} again
\end{itemize}

\begin{tcolorbox}[breakable,colback=teal!4,colframe=teal!35!black,title={Before you move on}]
What does \te{ఆకారం} mean? (\textquotedblleft{}A shape\textquotedblright{}.) Say it once more.
\end{tcolorbox}
\section[\texorpdfstring{\te{ఎందుకైనా} (endukainā)}{ఎందుకైనా (endukainā)}]{\te{ఎందుకైనా} (endukainā) — in case, to be on the safe side}

\label{lesson:TE-C295-endukaina}



\begin{quote}
Before the new one: say the Telugu for success, then the Telugu for a shape.
\end{quote}

\subsection*{You'll want to know: \te{ఎందుకైనా}}
\textbf{\te{ఎందుకైనా}} — \emph{endukainā} — \textquotedblleft{}in case, to be on the safe side\textquotedblright{}.

\begin{grammarlens}[title={What you've built}]
One more word for this chapter.
\end{grammarlens}

\subsection*{Your turn}
\begin{itemize}
\raggedright
  \item \emph{Say it:} \emph{endukainā}
  \item \emph{Say it:} \emph{endukainā}, once more
  \item \emph{Say it:} say \emph{endukainā}
  \item \emph{Recall:} say the Telugu for success, then the Telugu for a shape, then say \emph{endukainā} again
\end{itemize}

\begin{tcolorbox}[breakable,colback=teal!4,colframe=teal!35!black,title={Before you move on}]
What does \te{ఎందుకైనా} mean? (\textquotedblleft{}In case, to be on the safe side\textquotedblright{}.) Say it once more.
\end{tcolorbox}
\section[\texorpdfstring{\te{ఎందుకో} (endukō)}{ఎందుకో (endukō)}]{\te{ఎందుకో} (endukō) — for some reason}

\label{lesson:TE-C295-enduko}



\begin{quote}
Before the new one: say the Telugu for a shape, then the Telugu for in case.
\end{quote}

\subsection*{You'll want to know: \te{ఎందుకో}}
\textbf{\te{ఎందుకో}} — \emph{endukō} — \textquotedblleft{}for some reason\textquotedblright{}.

\begin{grammarlens}[title={What you've built}]
One more word for this chapter.
\end{grammarlens}

\subsection*{Your turn}
\begin{itemize}
\raggedright
  \item \emph{Say it:} \emph{endukō}
  \item \emph{Say it:} \emph{endukō}, once more
  \item \emph{Say it:} say \emph{endukō}
  \item \emph{Recall:} say the Telugu for a shape, then the Telugu for in case, then say \emph{endukō} again
\end{itemize}

\begin{tcolorbox}[breakable,colback=teal!4,colframe=teal!35!black,title={Before you move on}]
What does \te{ఎందుకో} mean? (\textquotedblleft{}For some reason\textquotedblright{}.) Say it once more.
\end{tcolorbox}
\section[\texorpdfstring{\te{ఎలాగైనా} (elāgainā)}{ఎలాగైనా (elāgainā)}]{\te{ఎలాగైనా} (elāgainā) — somehow, by any means}

\label{lesson:TE-C295-elagaina}



\begin{quote}
Before the new one: say the Telugu for in case, then the Telugu for for some reason.
\end{quote}

\subsection*{You'll want to know: \te{ఎలాగైనా}}
\textbf{\te{ఎలాగైనా}} — \emph{elāgainā} — \textquotedblleft{}somehow, by any means\textquotedblright{}.

\begin{grammarlens}[title={What you've built}]
One more word for this chapter.
\end{grammarlens}

\subsection*{Your turn}
\begin{itemize}
\raggedright
  \item \emph{Say it:} \emph{elāgainā}
  \item \emph{Say it:} \emph{elāgainā}, once more
  \item \emph{Say it:} say \emph{elāgainā}
  \item \emph{Recall:} say the Telugu for in case, then the Telugu for for some reason, then say \emph{elāgainā} again
\end{itemize}

\begin{tcolorbox}[breakable,colback=teal!4,colframe=teal!35!black,title={Before you move on}]
What does \te{ఎలాగైనా} mean? (\textquotedblleft{}Somehow, by any means\textquotedblright{}.) Say it once more.
\end{tcolorbox}
\section[(dialogue)]{Words from chapters 287-290 — a first pass}

\label{lesson:TE-R295-first-pass-words-from-chapters-287-290}



\begin{quote}
Say the words for for some reason and somehow.
\end{quote}

\subsection*{Your turn}
Say these aloud:

\begin{itemize}
\raggedright
  \item the chest, an accident, an ambulance, insurance, an emergency
  \item an operation, indigestion, constipation, diabetes, blood pressure
  \item Bhogi, a festival, a village fair, a celebration, sunrise
  \item daytime, a half of the day, a decade, a century, childhood
\end{itemize}

\begin{tcolorbox}[breakable,colback=teal!4,colframe=teal!35!black,title={Before you move on}]
The chest? (\textbf{\te{ఛాతీ}}, \emph{chātī}.) An accident? (\textbf{\te{ప్రమాదం}}, \emph{pramādaṁ}.) An ambulance? (\textbf{\te{అంబులెన్సు}}, \emph{ambulensu}.) Insurance? (\textbf{\te{బీమా}}, \emph{bīmā}.) An emergency? (\textbf{\te{అత్యవసరం}}, \emph{atyavasaraṁ}.) An operation? (\textbf{\te{ఆపరేషను}}, \emph{āparēṣanu}.) Indigestion? (\textbf{\te{అజీర్ణం}}, \emph{ajīrṇaṁ}.) Constipation? (\textbf{\te{మలబద్ధకం}}, \emph{malabaddhakaṁ}.) Diabetes? (\textbf{\te{మధుమేహం}}, \emph{madhumēhaṁ}.) Blood pressure? (\textbf{\te{రక్తపోటు}}, \emph{raktapōṭu}.) Bhogi? (\textbf{\te{భోగి}}, \emph{bhōgi}.) A festival? (\textbf{\te{పండుగ}}, \emph{paṇḍuga}.) A village fair? (\textbf{\te{జాతర}}, \emph{jātara}.) A celebration? (\textbf{\te{వేడుక}}, \emph{vēḍuka}.) Sunrise? (\textbf{\te{సూర్యోదయం}}, \emph{sūryōdayaṁ}.) Daytime? (\textbf{\te{పగలు}}, \emph{pagalu}.) A half of the day? (\textbf{\te{పూట}}, \emph{pūṭa}.) A decade? (\textbf{\te{దశాబ్దం}}, \emph{daśābdaṁ}.) A century? (\textbf{\te{శతాబ్దం}}, \emph{śatābdaṁ}.) Childhood? (\textbf{\te{బాల్యం}}, \emph{bālyaṁ}.)
\end{tcolorbox}
\section[(dialogue)]{Words from chapters 291-295 — a second pass}

\label{lesson:TE-R295-second-pass-words-from-chapters-291-295}



\begin{quote}
Say the words for south and a place.
\end{quote}

\subsection*{Your turn}
Say these aloud:

\begin{itemize}
\raggedright
  \item south, a place, a border, opposite, around
  \item a hut, a canal, a building, a fort, a factory
  \item a jail, a mosque, a church, the village council, a loft
  \item advice, moreover, by the way, almost, all in all
  \item success, a shape, in case, for some reason, somehow
\end{itemize}

\begin{tcolorbox}[breakable,colback=teal!4,colframe=teal!35!black,title={Before you move on}]
South? (\textbf{\te{దక్షిణం}}, \emph{dakṣiṇaṁ}.) A place? (\textbf{\te{ప్రదేశం}}, \emph{pradēśaṁ}.) A border? (\textbf{\te{సరిహద్దు}}, \emph{sarihaddu}.) Opposite? (\textbf{\te{ఎదురుగా}}, \emph{edurugā}.) Around? (\textbf{\te{చుట్టూ}}, \emph{cuṭṭū}.) A hut? (\textbf{\te{గుడిసె}}, \emph{guḍise}.) A canal? (\textbf{\te{కాలువ}}, \emph{kāluva}.) A building? (\textbf{\te{భవనం}}, \emph{bhavanaṁ}.) A fort? (\textbf{\te{కోట}}, \emph{kōṭa}.) A factory? (\textbf{\te{కార్ఖానా}}, \emph{kārkhānā}.) A jail? (\textbf{\te{జైలు}}, \emph{jailu}.) A mosque? (\textbf{\te{మసీదు}}, \emph{masīdu}.) A church? (\textbf{\te{చర్చి}}, \emph{carci}.) The village council? (\textbf{\te{పంచాయతీ}}, \emph{pañcāyatī}.) A loft? (\textbf{\te{అటక}}, \emph{aṭaka}.) Advice? (\textbf{\te{సలహా}}, \emph{salahā}.) Moreover? (\textbf{\te{పైగా}}, \emph{paigā}.) By the way? (\textbf{\te{అన్నట్టు}}, \emph{annaṭṭu}.) Almost? (\textbf{\te{దాదాపు}}, \emph{dādāpu}.) All in all? (\textbf{\te{మొత్తానికి}}, \emph{mottāniki}.) Success? (\textbf{\te{విజయం}}, \emph{vijayaṁ}.) A shape? (\textbf{\te{ఆకారం}}, \emph{ākāraṁ}.) In case? (\textbf{\te{ఎందుకైనా}}, \emph{endukainā}.) For some reason? (\textbf{\te{ఎందుకో}}, \emph{endukō}.) Somehow? (\textbf{\te{ఎలాగైనా}}, \emph{elāgainā}.)
\end{tcolorbox}
